\section{Future Work}
\label{sec6.3: future work}

The challenges identified in \ref{sec6.2: limitations} point toward several avenues for future research:

\begin{itemize}
    \item \textbf{Reducing Prediction Set Sizes Directly:} Beyond the Gram type filter we already used, we can use a second
          coarser filtering stage, for example something based on phage morphology or a nearest neighbor prefilter could narrow 
          the host level candidate space even further before the conformal envelope test is applied. 

    \item \textbf{Pooling for Low Sample Hosts:} Rather than calibrating all of the $54$ host envelopes 
           independently, a hierarchical shrinkage scheme could pool information across related hosts. Shrinking quantile 
           estimates for data scarce genera ($n_h < 10$) toward family or Gram level baselines would stabilize coverage 
           without requiring additional data.

    \item \textbf{Dimensionality Reduction:} To resolve high dimensional sparsity in multi 
           dimensional geometries ($K_h \approx 40$), future work should explore learning host specific linear or non linear 
           projections $\Phi_h: \mathbb{R}^{K_h} \to \mathbb{R}^{d}$ ($d \in [2, 5]$) designed to maximize infector/non 
           infector separation. 
        %    Learning a Mahalanobis distance metric may help us preserve local support for the Radial envelope in higher 
        %    dimensions and reduce reliance on the marginal fallback.

    \item \textbf{Specificity Control:} The impostor calibration experiment (\ref{sec5.3.7: host ablations}) shows 
          that naively incorporating non infector scores collapses the prediction sets to near total abstention. Future work 
          may investigate two sided nonconformity functions that jointly penalize false negatives and false positives, 
          this may bound the false discovery rate and yet avoid the coverage loss.
    
    \item \textbf{Handling of Missing Scores:} Instead of defaulting to the conservative marginal quantile when 
          functional categories are absent, future envelope constructions could condition fallbacks on the observed 
          missingness patterns. This may be done by observing the sparsity pattern, which itself may be carrying some 
          information about the host.

      \item \textbf{Dimensionality Reduction} As demonstrated by our core functional analysis (Section~\ref{sec5.3.6: 
           sufficiency}), most discriminative information lives in a low dimensional subspace. In higher dimensions ($K_h 
           \approx 40$), sample sparsity forces directional cones and grid bins to fall back onto conservative marginal 
           bounds. Future work should explore linear dimensionality reduction such as class specific Principal Component 
           Analysis (PCA) to extract dominant functional variance directions. Supervised linear/non-linear projections 
           $\Phi_h: \mathbb{R}^{K_h} \to \mathbb{R}^d$ ($d \in [2, 5]$) optimized to maximize infector/non infector 
           separation. Calibrating Radial and Strip envelopes within these compact, orthogonal PCA subspaces would 
           eliminate the curse of dimensionality while preserving directional sensitivity
\end{itemize}









